\documentclass[12pt,a4paper]{article}
\usepackage[utf8]{inputenc}
\usepackage[italian]{babel}
\usepackage{amsmath, amsfonts, amssymb}
\usepackage{graphicx}
\usepackage{float}
\usepackage{hyperref}
\usepackage{geometry}
\geometry{a4paper, margin=1in}
\usepackage{setspace}
\onehalfspacing

\begin{document}

% ==========================================
% PAGINA 1: FRONTESPIZIO (Stile Tesi)
% ==========================================
\begin{titlepage}
    \centering
    \vspace*{2cm}
    
    {\huge\bfseries Modellazione e Scoperta Parametrica della Propagazione Epidemica\par}
    \vspace{0.8cm}
    {\Large\itshape Metodi Numerici Classici vs Physics-Informed Neural Networks (PINN)\par}
    
    \vspace{3cm}
    {\Large Studente: \\ \textbf{Emanuele Barbagallo}\par}
    \vspace{1.5cm}
    {\large Esame di: \\ \textit{Numerical Methods for Scientific Computing}\par}
    
    \vfill
    
    %{\large Anno Accademico 2025/2026\par}
    %\vspace{1cm}
    {\large Appello del 21 Settembre 2026\par}
\end{titlepage}

% ==========================================
% PAGINA 2: ABSTRACT
% ==========================================
\newpage
\thispagestyle{empty}
\vspace*{4cm}
\begin{center}
    \textbf{\huge Abstract}
\end{center}
\vspace{1cm}
Il presente progetto affronta l'estrazione del coefficiente di diffusione spaziale ($D$) dai dati storici dell'epidemia COVID-19. Per risolvere questo problema inverso, vengono sviluppate e confrontate due architetture antitetiche: un metodo numerico classico (Differenze Finite e ottimizzazione di Levenberg-Marquardt) e un framework di Deep Learning all'avanguardia basato su \textit{Physics-Informed Neural Networks} (PINN). I risultati validano empiricamente la convergenza deterministica di entrambi gli approcci, portando alla luce i profondi trade-off esistenti tra l'efficienza algoritmica dei sistemi sparsi e la resilienza architetturale delle formulazioni \textit{mesh-free}.


% ==========================================
% PAGINA 3: INDICE
% ==========================================
\newpage
\tableofcontents
\newpage

% ==========================================
% CORPO DELLA RELAZIONE
% ==========================================

\section{Introduzione e Acquisizione dei Dati}
La modellazione matematica delle epidemie si basa storicamente su sistemi dinamici accoppiati. Tuttavia, per comprendere l'evoluzione puramente spaziale e territoriale del contagio, è necessario spostarsi nel campo delle Equazioni alle Derivate Parziali (PDE). 

Per alimentare il modello matematico con informazioni empiriche, è stato sviluppato uno script in Python (tramite libreria \texttt{pandas}) dedito all'acquisizione e alla pulizia dei dati. La scelta architetturale è stata quella di non archiviare file \texttt{.csv} statici di grandi dimensioni all'interno del progetto, ma di effettuare un download dinamico puntando direttamente alla repository GitHub ufficiale del Dipartimento della Protezione Civile. Questa scelta previene l'obsolescenza dei dati e rende l'intero progetto portatile.

\subsection{Filtraggio Temporale e Pulizia del Dataset}
Durante la manipolazione del \textit{dataframe}, sono stati affrontati diversi problemi intrinseci alla struttura dei dati governativi:
\begin{itemize}
    \item \textbf{Correzione del Timezone:} I dataset includevano originariamente il fuso orario UTC. Al fine di evitare collisioni critiche con i filtri temporali di Pandas (basati su formati "naive"), è stato operato un distacco del timezone (\texttt{tz\_localize(None)}), prevenendo errori algoritmici legati ai cambi di ora legale.
    \item \textbf{Isolamento della "Prima Ondata":} È stato operato un ritaglio temporale chirurgico, considerando esclusivamente l'intervallo dal 24 Febbraio 2020 al 31 Maggio 2020. Questa scelta non è casuale: isolare la primissima ondata garantisce di operare su un'epidemia "pura", priva di nuove varianti del virus, assenza di vaccini e interventi governativi territoriali asimmetrici (lockdown mirati localmente), i quali avrebbero introdotto fortissime non-linearità invalidando la stabilità formale del modello differenziale.
    \item \textbf{L'Anomalia del Golfo di Guinea:} Durante le fasi più concitate dell'emergenza, i dati "in fase di definizione" o privi di geolocalizzazione venivano registrati dalla Protezione Civile assegnando d'ufficio coordinate Latitudine = 0 e Longitudine = 0. In assenza di un'epurazione rigorosa di queste righe, il modello differenziale avrebbe collocato decine di migliaia di infetti nel bel mezzo dell'Oceano Atlantico (Golfo di Guinea), corrompendo irrimediabilmente tutta l'algebra spaziale del progetto. Tali coordinate anomale sono state sistematicamente espunte tramite filtraggio vettoriale.
\end{itemize}

\subsection{Mappatura Spaziale (Grid Mapping)}
L'ultimo passaggio dell'acquisizione dati consiste nella traduzione del sistema geografico. Le coordinate continue (Latitudine e Longitudine) delle province italiane non possono essere date in pasto direttamente a un metodo alle Differenze Finite. Per questo motivo, l'intero territorio nazionale è stato scalato e normalizzato all'interno di un dominio matematico bidimensionale unitario $[0, 1] \times [0, 1]$. In seguito, questo spazio continuo è stato discretizzato trasformandolo in una vera e propria "scacchiera", una griglia $50 \times 50$. Ogni provincia, in base alle sue coordinate, è stata agganciata in approssimazione al nodo più vicino di tale matrice, fornendo così la base topologica di calcolo esatta (i \textit{Data Points}) per le successive discretizzazioni del Laplaciano spaziale.


\section{L'Approccio Numerico Classico}

\subsection{L'Equazione alle Derivate Parziali (PDE)}
Il fenomeno epidemiologico spaziale è stato modellato attraverso l'Equazione di Diffusione, un'equazione parabolica la cui forma canonica è:
\begin{equation}
    \frac{\partial u}{\partial t} = D \nabla^2 u
\end{equation}
dove $u(x,y,t)$ rappresenta la densità spaziotemporale dei contagi e $D$ rappresenta il coefficiente di diffusione, la nostra incognita macroscopica. L'operatore differenziale $\nabla^2$ è il Laplaciano spaziale bidimensionale, definito analiticamente come la somma delle derivate parziali seconde:
\begin{equation}
    \nabla^2 u = \frac{\partial^2 u}{\partial x^2} + \frac{\partial^2 u}{\partial y^2}
\end{equation}
Fisicamente, la derivata seconda misura la concavità topologica: determina matematicamente se in una specifica provincia vi è un accumulo anomalo di contagi rispetto alle adiacenti (che implicherebbe un flusso in ingresso) o una dispersione verso l'esterno.

\subsection{Sviluppo di Taylor e Stencil a 5 Punti}
Poiché i calcolatori elaborano esclusivamente domini discreti, è imperativo convertire le derivate continue in algebra lineare. L'analisi numerica affronta questo scoglio tramite lo \textbf{Sviluppo in Serie di Taylor}. Valutando i punti adiacenti $x_{i+1}$ e $x_{i-1}$ e troncando la serie al secondo ordine, si ottiene la celebre approssimazione alle Differenze Finite Centrali. Sull'asse orizzontale si ha:
\begin{equation}
    \frac{\partial^2 u}{\partial x^2} \approx \frac{u_{i+1, j} - 2u_{i,j} + u_{i-1, j}}{\Delta x^2}
\end{equation}
Replicando la medesima operazione sull'asse verticale (Nord-Sud) per $\Delta y^2$ e assumendo che la griglia territoriale (la nostra discretizzazione matriciale 50x50) presenti un campionamento perfettamente quadrato ($\Delta x = \Delta y = h$), le due espressioni possono essere fuse per ottenere l'approssimazione integrale del Laplaciano spaziale:
\begin{equation}
    \nabla^2 u \approx \frac{u_{i+1, j} + u_{i-1, j} + u_{i, j+1} + u_{i, j-1} - 4u_{i,j}}{h^2}
\end{equation}
Da questa deduzione analitica emerge la topologia della matrice di calcolo: il nodo centrale contribuisce con un coefficiente di $-4$ (costituendo la diagonale principale della matrice di rigidezza), mentre i quattro vicini cardinali (Nord, Sud, Est, Ovest) contribuiscono con coefficienti $+1$ (costituendo le diagonali secondarie). Questa struttura geometrica "a croce" è nota nella letteratura come \textbf{Stencil a 5 punti}.

Applicando lo Stencil all'intero dominio geografico vettorializzato, la PDE continua si trasforma rigorosamente in un gigantesco sistema di equazioni algebriche lineari della forma $A \mathbf{u} = \mathbf{b}$, dove $A$ è la matrice dei coefficienti. 

\subsubsection{La Maledizione della Dimensionalità e le Matrici Sparse}
Un aspetto critico dell'approccio numerico è la dipendenza dalla densità della griglia. Qualora definissimo una griglia molto fitta (es. $1000 \times 1000$), il vettore di stato $\mathbf{u}$ avrebbe 1 milione di incognite, e la matrice $A$ avrebbe dimensioni $10^6 \times 10^6$. Salvare una matrice densa di tale entità saturerebbe istantaneamente la RAM.
Per eludere questa limitazione fisica, la matrice $A$ è stata codificata sfruttando le strutture dati per le matrici sparse di \texttt{scipy.sparse}, specificatamente nel formato \textit{Compressed Sparse Row} (CSR). Questo formato memorizza esclusivamente gli elementi non nulli, abbattendo drasticamente la complessità spaziale e permettendo il calcolo del Laplaciano in tempi ridottissimi.

\subsection{Solutori Iterativi e Sottospazi di Krylov (Forward Problem)}
Calcolare l'evoluzione temporale dei contagi assumendo noto il parametro $D$ costituisce il "Forward Problem".
A causa delle enormi dimensioni della matrice sparsa descritta precedentemente, metodi di fattorizzazione diretta (come l'eliminazione gaussiana o la decomposizione LU) risultano proibitivi a causa del fenomeno del \textit{fill-in} (la creazione di nuovi elementi non nulli che distrugge la sparsità).

Per ovviare a questo problema, la risoluzione del sistema $A \mathbf{x} = \mathbf{b}$ è stata affidata a metodi iterativi basati sulle proiezioni nei \textbf{Sottospazi di Krylov}, definiti matematicamente come:
\begin{equation}
    \mathcal{K}_m(A, r_0) = \text{span} \{ r_0, A r_0, A^2 r_0, \dots, A^{m-1} r_0 \}
\end{equation}
dove $r_0 = \mathbf{b} - A \mathbf{x}_0$ è il residuo iniziale. La potenza computazionale di questi metodi deriva dall'esecuzione esclusiva di moltiplicazioni matrice-vettore. In questo contesto, il formato CSR adottato nella fase di discretizzazione si rivela imbattibile: saltando a livello hardware le coordinate degli zeri, la CPU evita decine di milioni di calcoli a vuoto ($0 \times x_i = 0$), abbattendo lo sforzo iterativo di interi ordini di grandezza.

All'interno dell'architettura sono stati messi a confronto due solutori principali:
\begin{itemize}
    \item \textbf{Gradiente Coniugato (CG):} Teoricamente ottimale in termini di memoria e velocità. Tuttavia, affinché il CG converga, la matrice $A$ deve essere rigorosamente \textbf{Simmetrica e Definita Positiva (SPD)}. La matrice del Laplaciano generata dallo stencil possiede coefficienti diagonali pari a $-4$, risultando definita negativa. Prima di avviare il risolutore, si è reso necessario moltiplicare l'intero sistema per $-1$, ottenendo una matrice SPD pura. Fisicamente, la \textit{simmetria} ($A = A^T$) riflette l'isotropia della diffusione virale, mentre la \textit{definitezza positiva} ($x^T A x > 0$) attesta che il sistema termodinamico sia dissipativo, impedendo la generazione di contagi spontanei.
    \item \textbf{GMRES (Generalized Minimal Residual Method):} Implementato come solutore di contingenza. Qualora il modello epidemiologico venisse esteso con termini convettivi asimmetrici, la matrice $A$ perderebbe la propria simmetria, causando il crash algoritmico del CG. Il GMRES minimizza la norma del residuo su matrici asimmetriche estraendo vettori dal sottospazio di Krylov. Il cuore matematico del GMRES è l'\textbf{Iterazione di Arnoldi}, un processo che costruisce, passo dopo passo, una base ortonormale per lo spazio di Krylov. Ad ogni iterazione, l'algoritmo proietta il grande sistema lineare su questo piccolo sottospazio, riducendolo a un problema ai minimi quadrati.

\end{itemize}

\subsection{Ottimizzazione Non Lineare: Levenberg-Marquardt (Inverse Problem)}
La vera sfida computazionale (l'obiettivo del progetto) risiede nell'Inverse Problem: noi possediamo i dati reali storici, ma ignoriamo il valore di $D$ che li ha generati. 
L'estrazione del parametro è stata formulata come un problema di Data Fitting ai minimi quadrati. L'obiettivo matematico è trovare il $D^*$ tale per cui:
\begin{equation}
    D^* = \arg\min_D \sum_{i} \left( u_{pred}(D, x_i) - u_{true}(x_i) \right)^2
\end{equation}

La risoluzione di questo sistema non lineare è stata affidata all'algoritmo di \textbf{Levenberg-Marquardt} (LM), formalmente noto come metodo dei Minimi Quadrati Smorzati (\textit{Damped Least-Squares}). L'algoritmo calcola l'aggiornamento del parametro ($\Delta D$) ad ogni iterazione risolvendo la seguente equazione algebrica:
\begin{equation}
    (J^T J + \lambda I) \Delta D = J^T r
\end{equation}
dove i singoli termini rappresentano:
\begin{itemize}
    \item $r = u_{true} - u_{pred}(D)$: è il vettore dei residui (la discrepanza tra i contagi reali e quelli previsti).
    \item $J$: è la matrice Jacobiana, contenente le derivate parziali dei residui rispetto al parametro $D$.
    \item $J^T J$: costituisce l'approssimazione definita positiva della matrice Hessiana (comportamento tipico del metodo di Gauss-Newton).
    \item $\lambda$: è lo scalare di smorzamento (\textit{damping parameter}) regolato dinamicamente.
    \item $I$: è la matrice Identità.
\end{itemize}
La genialità matematica di LM risiede nella sua architettura ibrida governata da $\lambda$. Quando la ricerca si trova lontana dal minimo globale ($\lambda$ elevato), il termine $\lambda I$ domina l'equazione, forzando l'algoritmo a comportarsi come un \textit{Gradient Descent} puro per evitare divergenze; man mano che si approccia alla valle del minimo ($\lambda \to 0$), il termine si annulla e LM muta organicamente nel metodo di \textit{Gauss-Newton}, sfruttando l'approssimazione dell'Hessiana ($J^T J$) per innescare una fulminea convergenza quadratica finale verso la soluzione esatta.


\subsubsection{Validazione Algoritmica e "Controfigura Matematica"}
L'addestramento di un ottimizzatore esige decine di valutazioni della funzione di costo. Se durante la stesura del codice l'algoritmo avesse dovuto generare ripetutamente la matrice sparsa del Laplaciano $2500 \times 2500$ e risolvere i sistemi di Krylov, ogni tentativo di \textit{debugging} avrebbe impiegato ore. 
Per aggirare questo ostacolo prestazionale, il solutore LM è stato validato in isolamento utilizzando una \textbf{"controfigura matematica"}: un modello surrogato ultra-leggero, come un banale decadimento esponenziale del tipo $u = u_0 \cdot e^{-2D}$. Fissando una \textit{ground truth} arbitraria per testarne il funzionamento (es. $D_{true} = 0.3500$) e fornendo all'algoritmo una supposizione di partenza completamente errata ($D_{guess} = 1.0$), Levenberg-Marquardt ha operato una discesa alla cieca nello spazio delle soluzioni, convergendo al millesimo verso il parametro target. 

Accertata inequivocabilmente l'infallibilità matematica dell'ottimizzatore, l'equazione surrogata è stata espunta dall'architettura e sostituita definitivamente dal vero motore differenziale (il Forward Problem tramite GMRES/CG). Questo elegante stratagemma di \textit{mocking} ha garantito la robustezza del codice finale.


\section{L'Approccio Deep Learning: Physics-Informed Neural Networks}
Come antitesi ingegneristica all'approccio classico, si è esplorata una soluzione basata sull'intelligenza artificiale, specificatamente le PINN. 
Una rete neurale tradizionale è intrinsecamente una \textit{Black Box}: interpola i dati tramite gradienti statistici, ignorando totalmente le leggi dell'universo fisico, rischiando di generare previsioni termodinamicamente impossibili.
La PINN sovverte questo paradigma codificando l'Equazione Differenziale Parziale direttamente all'interno della \textit{Loss Function} (funzione di costo). La rete deve minimizzare due componenti simultaneamente:
\begin{equation}
    \mathcal{L}_{total} = \lambda_1 \mathcal{L}_{data} + \lambda_2 \mathcal{L}_{physics}
\end{equation}
dove $\mathcal{L}_{data}$ misura l'errore rispetto ai dati dei contagi reali (come in una rete classica), mentre $\mathcal{L}_{physics}$ misura di quanto l'output della rete violi la forma matematica $\frac{\partial u}{\partial t} - D \nabla^2 u = 0$. In questo processo, il calcolo delle derivate spaziali e temporali non è soggetto all'errore di troncamento delle griglie, bensì calcolato in modo analitico esatto grazie al motore di \textit{Automatic Differentiation} (Autograd) di librerie come PyTorch.

\subsection{L'Architettura Mesh-Free e la Resilienza Topologica}
Il trionfo concettuale delle PINN risiede nella loro natura indipendente dalle griglie (\textit{Mesh-Free}). Mentre i metodi alle differenze finite pretendono una discretizzazione spaziale rigida ed equispaziata, la rete neurale astrae il concetto di spazio campionando in maniera continua un dominio vettoriale, avvalendosi di due distinte tipologie di coordinate per l'addestramento:
\begin{itemize}
    \item \textbf{Data Points:} Rappresentano le coordinate geospaziali in cui disponiamo delle misurazioni storiche reali dei contagi (i tamponi accertati della Protezione Civile). Su questi punti viene calcolata esclusivamente la $\mathcal{L}_{data}$, costringendo la rete ad ancorare fermamente le proprie predizioni alla realtà clinica empirica.
    \item \textbf{Collocation Points:} Costituiscono il vero nucleo innovativo dell'architettura informata dalla fisica. Si tratta di punti vuoti generati (o "spruzzati") in maniera puramente stocastica su tutto il dominio spaziotemporale, comprese province remote, laghi o aree prive di qualsiasi sensore di rilevamento epidemiologico. Su questi nodi fittizi non possediamo \textit{ground truth}, ma imponiamo alla rete di rispettare pedissequamente la $\mathcal{L}_{physics}$. In sostanza, l'algoritmo impara a propagare il contagio differenziale in modo coerente anche laddove non ha mai "osservato" un dato.
\end{itemize}
Questa struttura computazionale bipartita garantisce una resilienza topologica impareggiabile: qualora venissero a mancare misurazioni in vaste porzioni di territorio, le matrici FDM presenterebbero lacune strutturali (divenendo singolari e mandando in blocco i solutori di Krylov). La PINN, al contrario, compensa istantaneamente l'assenza di Data Points inondando le "zone d'ombra" con migliaia di Collocation Points, deducendo il tracciato epidemiologico in virtù della sola coerenza termodinamica.


\subsection{La Gradient Pathology: Adam vs L-BFGS}
Durante la fase di addestramento (Parameter Discovery neurale), la rete cerca di calibrare il parametro $D$ associato come \texttt{nn.Parameter} addestrabile. In questa fase è emerso un noto problema della letteratura PINN: la \textit{Gradient Pathology}.
I gradienti derivanti dalla \textit{Data Loss} e quelli derivanti dalla \textit{Physics Loss} possiedono magnitudo sbilanciate o spingono verso direzioni opposte, alterando la curvatura geometrica dello spazio di ottimizzazione.

Di fronte a questo scenario:
\begin{itemize}
    \item \textbf{Adam:} L'ottimizzatore stocastico Adam (basato esclusivamente sulle derivate del primo ordine e sull'inerzia) rimane intrappolato nei minimi locali generati dall'anomalia. Empiricamente, si osserva uno stallo perpetuo dell'errore (l'estrazione del parametro si ferma a distanza dalla realtà, generando una curva piatta di fallimento).
    \item \textbf{L-BFGS:} Per aggirare l'ostacolo, il protocollo di addestramento è stato commutato su L-BFGS, un ottimizzatore \textit{Quasi-Newton} del secondo ordine. Approssimando l'inversa della matrice Hessiana, L-BFGS riesce a decodificare la complessità della curvatura anomala, abbattendo verticalmente la Loss function e convergendo trionfalmente al vero minimo globale, estraendo con precisione il valore deterministico esatto.
\end{itemize}

\section{Conclusioni: Trade-off Prestazionali e Architetturali}
Il traguardo dell'elaborato consisteva nell'operare un confronto diretto. 
Sotto il profilo dell'accuratezza formale, la ricerca si chiude in un netto pareggio tecnico: sia il Metodo Numerico (con Levenberg-Marquardt) sia la Physics-Informed Neural Network (con L-BFGS) hanno individuato con precisione microscopica il vero coefficiente di diffusione ($D = 0.3500$). 
Questo risultato garantisce la validità incrociata di entrambi i framework.

Ciononostante, il confronto solleva profondi trade-off di natura architetturale, sui quali ruota la valutazione finale del progetto:
\begin{enumerate}
    \item \textbf{Costi Computazionali ed Efficienza:} Se il contesto applicativo verte sull'efficienza di calcolo in fase di discovery, l'approccio numerico distrugge ogni competizione. La solidità dei solutori iterativi e l'impiego delle matrici sparse (CSR) permettono all'algebra lineare classica di risolvere l'Inverse Problem in circa \textbf{0.05 secondi}. Di contro, le reti neurali necessitano di migliaia di epoche, richiedendo l'uso di acceleratori hardware (Cluster GPU) e tempistiche che variano dai svariati minuti alle ore.
    \item \textbf{Il Costo dell'Inferenza:} Se i metodi numerici impongono un tempo quasi nullo di calibrazione, possiedono però un costo di inferenza iterativo. Per simulare 1000 giorni futuri, l'algoritmo classico deve ri-eseguire il solutore per 1000 passi temporali sequenziali. La PINN sconta il suo prezzo all'inizio (\textit{Upfront Cost} di training), ma una volta calibrata, offre un costo di predizione istantaneo $\mathcal{O}(1)$ (basta fornire alla rete il tempo $t=1000$ per ottenere la matrice).
    \item \textbf{Dimensionalità e Resilienza Topologica:} I metodi alle Differenze Finite soffrono terribilmente della maledizione della dimensionalità (esplosione delle matrici al crescere dello spazio in 3D o 4D) e crollano di fronte a campionamenti non equispaziati. La PINN, grazie all'emancipazione dalla griglia (Mesh-Free) e alla differenziazione automatica, garantisce un'estensibilità dimensionale pressoché gratuita e un'immunità formidabile alla frammentazione spaziale.
\end{enumerate}

In conclusione, la trattazione dimostra che l'impiego dei metodi numerici iterativi classici debba essere la scelta dominante in scenari controllati, caratterizzati da griglie regolari e geometrie ben definite. Laddove però la fluidodinamica o l'epidemiologia ci pongano dinanzi ad assenza cronica di dati, a sensori caotici, o a modelli spaziali ad alta complessità dimensionale, l'integrazione delle leggi di conservazione all'interno dei framework di Intelligenza Artificiale (PINN) non rappresenta solo un'alternativa intrigante, ma la nuova frontiera obbligata della moderna analisi numerica.

\end{document}
